% `tex.web` §336's third branch: an `\outer` control sequence in *skipped*
% conditional text.
%
%     print_err("Incomplete "); print_cmd_chr(if_test,cur_if);
%     print("; all text was ignored after line "); print_int(skip_line);
%     help3("A forbidden control sequence occurred in skipped text.")
%       ("This kind of error happens when you say `\if...' and forget")
%       ("the matching `\fi'. I've inserted a `\fi'; this might work.");
%     if cur_cs<>0 then cur_cs:=0
%     else help_line[2]:="The file ended while I was skipping conditional text.";
%     cur_tok:=cs_token_flag+frozen_fi; ins_error;
%
% Skipped text is still scanned, so the same two things that end a balanced
% scan end a skip — and here the recovery is a `\fi` rather than a `}`, with
% §495's `cur_if` and `skip_line` naming which conditional it belongs to and
% where its text began.
%
% The consequence to get right is what the inserted `\fi` finds. §336 leaves
% the conditional on `cond_ptr` and lets that `\fi` close it, so a skip that had
% not yet counted its conditional — a false branch, or an `\ifcase` looking for
% its case — must count it now. Otherwise the `\fi` arrives to find nothing
% open and there is an `Extra \fi` TeX does not have.
%
% The errors after each recovery are TeX's own and are the point of the case:
% the `\else` and `\fi` that were meant to close the conditional are now extra,
% because the inserted `\fi` closed it early.
\catcode`\{=1 \catcode`\}=2 \catcode`\#=6
\outer\def\bad{X}

\message{PROBE:one=before}
\iffalse
\bad
\else
\message{PROBE:two=else}
\fi
\message{PROBE:three=after}

% ...and in an `\ifcase` looking for its case, which counts its conditional the
% same way and names itself differently.
\ifcase 2 a\or b\bad\or c\fi
\message{PROBE:four=done}
\end
